\documentclass[10pt,a4paper]{amsart}
\usepackage[margin=19mm]{geometry}
\usepackage{amsmath,amssymb,mathtools}
\usepackage[scr=rsfs,cal=euler]{mathalfa}
\usepackage{xfrac}
\usepackage[unicode,hidelinks,bookmarks=false]{hyperref}
\newcommand{\ie}{i.e.\ }
\newcommand{\cf}{cf.\ }
\newcommand{\resp}{respectively\ }
\newcommand{\Subsection}{\subsection}
\providecommand{\nobreakdash}{\nobreak\hbox{-}}
\setlength{\emergencystretch}{2em}
\multlinegap0pt
\usepackage{bm}
\usepackage{bbm}
\usepackage{dsfont}
\usepackage{xspace}
\usepackage{cancel}
\usepackage{mathtools}
\usepackage{amsthm}
\makeatletter
\@ifundefined{theorem}{\newtheorem{theorem}{Theorem}}{}
\makeatother
\title[Dynamics of rotationally invariant polynomial root sets unde]{Dynamics of rotationally invariant polynomial root sets under iterated differentiations}
\author{Andr\'e Galligo, Joseph Najnudel, Truong Vu}
\date{Source: arXiv:2506.06263v2 (CC BY 4.0)}
\begin{document}
\maketitle

\noindent\emph{Source and licence.} Dynamics of rotationally invariant polynomial root sets under iterated differentiations. Version arXiv:2506.06263v2 (CC BY 4.0), distributed under the Creative Commons Attribution 4.0 International licence, as recorded in the arXiv licence metadata for this exact version and shown on the abstract page https://arxiv.org/abs/2506.06263v2. Full rights notice: licenses/arxiv-2506.06263-CC-BY-4.0.txt.\par\smallskip

\noindent\textit{Editorial scope.} Counted unit: the Theorem whose source label is \texttt{monotonicity} in the article (source0.tex lines 202--220, source label \texttt{monotonicity}), delivered together with the article's own complete original proof of it (source0.tex lines 221--328, one \texttt{proof} environment of 108 lines). No mathematics is written, repaired or re-derived: statement and proof are the article's own text line for line. Required context retained uncounted: none. Every definition, display and label of those lines is retained unchanged. Omitted from this package and not redistributed: 1--201, 329--1416. Nothing inside a retained excerpt is omitted or altered: every retained body is delivered exactly as the article's lines are written, so no omission receipt of any kind accompanies this package. Citations: the retained excerpts carry no citation command at all, so no bibliography entry of the article is transcribed and no reference is redistributed. Reference adaptation: none was needed and none was made; every reference inside the delivered unit resolves inside it, so no unresolved reference or citation is produced. Printed slips kept literal: no printed word, symbol, hypothesis or number of the counted statement or of its proof was altered. Layout and class: the article's own class is \texttt{amsart} with packages including geometry, graphicx, comment, float, amssymb, amsmath, epsfig, subfigure, epstopdf, enumerate, appendix, xcolor, hyperref; the wrapper is the lane's pinned \texttt{amsart} editorial wrapper at 10pt on A4, which restates the theorem-like environments the retained text needs and adds the article's own macro definitions that the retained text needs (none), with extra wrapper packages (bm, bbm, dsfont, xspace, cancel, mathtools). The delivered numbering is the unit's own. Assets: no retained span contains a figure environment, a graphics-inclusion command, a PDF-page inclusion, a table or a TikZ picture; no image file is copied into this package, the package declares no asset dependency and no image file is redistributed with it. Licence: version arXiv:2506.06263v2 of this article is distributed under the Creative Commons Attribution 4.0 International licence, as recorded in the arXiv licence metadata for this exact version and shown on the abstract page https://arxiv.org/abs/2506.06263v2. A full rights notice is delivered at licenses/arxiv-2506.06263-CC-BY-4.0.txt. Characters: every retained excerpt is pure ASCII, so the delivered engine, XeLaTeX with its default Latin Modern text font, reports no missing character.
\begin{theorem} \label{monotonicity}
Let $n \geq 2$, $w_1, \dots, w_n > 0$. Then, the function from $$ \Delta_n := \{ (z_1, \dots, z_n)\in  \mathbb{R}^n, z_1 \leq z_2 \leq \dots \leq z_n \} $$ to $\Delta_{n-1}$
such that the image of $(z_1, \dots, z_n)$ 
is the nondecreasing sequence of roots of the polynomial 
$$\sum_{j=1}^n w_j \prod_{1 \leq \ell \leq n, \ell \neq j} (z - z_\ell) \in \mathbb{R}_{n-1}(z),$$
counted with multiplicity, is increasing for the partial order $\preccurlyeq$ given by the following definition: on $\Delta_p$, 
$$(z_1, \dots, z_p) \preccurlyeq (z'_1, \dots, z'_p) $$
if and only if $z_j \leq z'_j$ for $1 \leq j \leq p$.

Moreover, if $w_1 = w_2 = \dots = w_n$, this function is $n/(n-1)$-Lipschitz for the L\'evy metric $L$ between empirical measures, which is defined for two probability measures $\mathbb{P}$ and 
$\mathbb{Q}$ by 
$$L(\mathbb{P}, \mathbb{Q}) = 
\inf \{\varepsilon > 0, \forall x \in \mathbb{R},
F_{\mathbb{P}} (x - \varepsilon ) - \varepsilon \leq F_{\mathbb{Q}}(x) \leq F_{\mathbb{P}} (x + \varepsilon ) + \varepsilon  \},$$
where $F_{\mathbb{P}}$ is the distribution function of $\mathbb{P}$ and $F_{\mathbb{Q}}$ is the distribution function of $\mathbb{Q}$. 



\end{theorem}

\begin{proof}
Let us assume 
$$(z_1, \dots, z_n) \preccurlyeq (z'_1, \dots, z'_n)$$
in $\Delta_n$. 
If we take into account multiplicities, for $1 \leq s \leq n-1$, the $s$-th smallest root of the two polynomials 
$$P : z \mapsto \sum_{j=1}^n w_j \prod_{1 \leq \ell \leq n, \ell \neq j} (z - z_\ell)$$ 
and 
$$Q : z \mapsto \sum_{j=1}^n w_j \prod_{1 \leq \ell \leq n, \ell \neq j} (z - z'_\ell)$$ 
lie in the intervals $[z_s, z_{s+1}]$ and $[z'_s, z'_{s+1}]$, respectively. If $z_s = z_{s+1}$ or 
$z'_s = z'_{s+1}$, we have $z_s \leq z_{s+1} \leq z'_s \leq z'_{s+1}$. Hence, the $s$-th root of $P$ is at most 
the $s$-th root of $Q$. We now assume that $z_s < z_{s+1}$
and $z'_s < z'_{s+1}$, and we denote by $\mu$ the $s$-th root of $P$. If $\mu \leq z'_s$, $\mu$ is at most the $s$-th root of $Q$.
Otherwise, $\mu$ is the unique root in $(z'_s, z_{s+1})$
of the rational function 
$$z \mapsto \sum_{j=1}^n \frac{w_j}{z - z_j}.$$
Since $z'_s < \mu < z_{s+1}$, $\mu$ is strictly at the right 
of all points in the interval $[z_r, z'_r]$ for $r \leq s$ and strictly 
at the left for $r \geq s+1$. We deduce that in both cases,
$1/(\mu - z)$ is well-defined and increasing in $z \in [z_r, z'_r]$. Hence, 
$$ \sum_{j = 1}^n \frac{w_j}{ \mu - z'_j} \geq
\sum_{j = 1}^n \frac{w_j}{ \mu - z_j}  = 0.$$
Now, the $s$-th root $\nu$ of $Q$ is the unique root
in $(z'_s , z'_{s+1})$ of the rational function
$$z \mapsto \sum_{j=1}^n \frac{w_j}{z - z'_j}.$$
Since this rational function is decreasing on $(z'_s , z'_{s+1})$ and is nonnegative at $\mu$, which is in this interval, 
we deduce that $\nu \geq \mu$. 
We have proven that the function in the proposition is nondecreasing. 
It is strictly increasing because a simple observation of the two leading coefficients shows that for 
 $$(z_1, \dots, z_n) \preccurlyeq (z'_1, \dots, z'_n)$$ 
 and 
 $$(z_1, \dots, z_n) \neq (z'_1, \dots, z'_n),$$
the sum of the roots of $Q$ is strictly larger than the sum of the roots of $P$. 

For the Lipschitz property in the case 
$w_1 = \dots = w_n$, let us assume that 
for $(z_1, \dots, z_n)$ and 
$(z'_1, \dots, z'_n)$ in $\Delta_n$, the corresponding empirical measures are at distance strictly smaller than $\varepsilon \in (0,1)$. 
In this case, for $1 \leq s \leq n$, applying the definition
of the L\'evy distance to $x < z_s - \varepsilon$,  and letting $x \rightarrow z_s - \varepsilon$, we deduce that the number of points 
among $z'_1, \dots, z'_n$ which are strictly smaller than 
$z_s - \varepsilon$ is at most $n ( (s-1)/n + \varepsilon)$, and then at most $s - 1 + \lfloor n \varepsilon \rfloor $. Hence, 
$z'_{s + \lfloor n \varepsilon \rfloor}  \geq z_s - \varepsilon$ as soon 
as $s + \lfloor n \varepsilon \rfloor \leq n$. 
We deduce that
$$(z'_1 - A, \dots, z'_{ \lfloor n \varepsilon \rfloor} - A, z_1 - \varepsilon, \dots, z_{n -\lfloor n \varepsilon \rfloor} - \varepsilon)
\preccurlyeq (z'_1, \dots, z'_n) $$
where $A > 0$ is sufficiently large, in order to have the left-hand side in $\Delta_n$. 
We then have the same inequality 
between the roots of the polynomials of degree $n-1$
constructed in the statement of the proposition. 
Let us compare the polynomials of degree $n-1$
constructed from 
$$(z'_1 - A, \dots, z'_{ \lfloor n \varepsilon \rfloor} - A, z_1 - \varepsilon, \dots, z_{n -\lfloor n \varepsilon \rfloor} - \varepsilon)$$ and from 
$$ (z_1 - \varepsilon, \dots, z_n - \varepsilon).$$
For $1 \leq r \leq n - 1- \lfloor n \varepsilon \rfloor$, 
the $(r + \lfloor n \varepsilon \rfloor)$-th 
root of the first polynomial, and the $r$-th root of the second polynomial are both in the interval $[z_r - \varepsilon, z_{r+1} - \varepsilon]$, when we take into account multiplicities. If $z_r < z_{r+1}$, these roots are the roots of rational functions, which are decreasing in $(z_r- \varepsilon, z_{r+1}- \varepsilon)$, 
the rational function constructed from 
$(z'_1 - A, \dots, z'_{ \lfloor n \varepsilon \rfloor} - A, z_1 - \varepsilon, \dots, z_{n -\lfloor n \varepsilon \rfloor} - \varepsilon)$ being larger than the rational function constructed from $(z_1 - \varepsilon, \dots, z_n - \varepsilon)$ at each point on $(z_r- \varepsilon, z_{r+1}- \varepsilon)$, because one goes from the first rational function to the second by replacing positive terms by negative terms, keeping the other terms unchanged: notice that here, we use the fact that $w_1 = w_2 = \dots = w_n$.  The $(r + \lfloor n \varepsilon \rfloor)$-th zero of the polynomial constructed 
from $(z'_1 - A, \dots, z'_{ \lfloor n \varepsilon \rfloor} - A, z_1 - \varepsilon, \dots, z_{n -\lfloor n \varepsilon \rfloor} - \varepsilon)$ is then at least equal to the 
$r$-th zero of the polynomial constructed from 
$(z_1- \varepsilon, \dots, z_n - \varepsilon)$. Hence, the 
$(r + \lfloor n \varepsilon \rfloor)$-th zero of the polynomial constructed from $(z'_1, \dots, z'_n)$ is at least the $r$-th zero $\mu_r$ of the 
polynomial constructed from $(z_1, \dots, z_n)$, minus $\varepsilon$. 
If $F$ and $G$ are the distribution functions of the empirical distribution of the roots of the polynomials of degree $n-1$ constructed from 
$(z_1, \dots, z_n)$ and $(z'_1, \dots, z'_n)$, we deduce that for $1 \leq r \leq n - 1 - \lfloor n \varepsilon \rfloor$,  
$x < \mu_r - \varepsilon$, 
and $x \geq  \mu_{r-1} - \varepsilon$ when $r \geq 2$, 
$$G(x) \leq \frac{1}{n-1} (r - 1 + \lfloor n \varepsilon \rfloor) \leq F(\mu_{r-1})   + \frac{n}{n-1} \varepsilon  \leq F(x + \varepsilon) +  \frac{n}{n-1} \varepsilon 
$$
when $r \geq 2$, and 
$$G(x) \leq \frac{1}{n-1} (\lfloor n \varepsilon \rfloor) \leq   \frac{n}{n-1} \varepsilon  \leq F(x + \varepsilon) +  \frac{n}{n-1} \varepsilon 
$$
when $r = 1$. 
We then get 
$$G(x) \leq  F(x + \varepsilon) +  \frac{n}{n-1} \varepsilon $$
for all $x < \mu_{n- 1 - \lfloor n \varepsilon \rfloor} - \varepsilon$, if $n- 1 - \lfloor n \varepsilon \rfloor \geq 1$. 
If $n - 1 -\lfloor n \varepsilon \rfloor \geq 1 $
and $x \geq \mu_{n- 1 - \lfloor n \varepsilon \rfloor}
- \varepsilon$, 
we get 
$$  F(x + \varepsilon) +  \frac{n}{n-1}  \varepsilon
\geq \frac{n - 1 -\lfloor n \varepsilon \rfloor}{n-1}
+ \frac{n}{n-1} \varepsilon \geq 1 \geq G(x),$$
and if $n - 1 - \lfloor n \varepsilon \rfloor = 0$, 
we have $n \varepsilon \geq n-1$ and then for all $x \in \mathbb{R}$, 
$$ F(x + \varepsilon) +  \frac{n}{n-1} \varepsilon
\geq 1 \geq G(x).$$
Hence, in any case, we have proven 
$$G(x) \leq F(x + \varepsilon) + \frac{n}{n-1}  \varepsilon.$$
Now, if we change all the points to their opposite and reverse their order, the distribution functions 
are changed via the map $F \mapsto \widetilde{F}$ where 
$$\widetilde{F} (x) = 1 - F((-x)-),$$
where $F((-x)-)$ denotes the left-limit of $F$ at $-x$,  
and the L\'evy distance between empirical measures does not change: indeed, increasing $\varepsilon$ by an arbitrarily small positive quantity in the definition of the distance absorbs possible errors due to the introduction of left limits of distribution functions. 
Applying the reasoning above after doing this transformation gives, for all $x \in \mathbb{R}$, 
with obvious notation, 
$$1 - G(x-) = \widetilde{G} (-x) \leq \widetilde{F} (-x + \varepsilon) + \frac{n}{n-1} \varepsilon
= 1 - F((x - \varepsilon)-) + \frac{n}{n-1} \varepsilon$$
and then
$$G(x) \geq G(x-) \geq F((x-\varepsilon)-) - \frac{n}{n-1} \varepsilon
\geq F(x-\varepsilon- \eta) - \frac{n}{n-1}
(\varepsilon + \eta),$$
for arbitrarily small $\eta > 0$. 
This inequality, combined with the upper bound above,
implies, after letting $\eta \rightarrow 0$, the Lipschitz property stated in the theorem. 

\end{proof}

\end{document}
